\clearpage
\fancyhead[L]{\small MERIDIANKREIS.}
\hypertarget{native-434-part-lower}{}
\begin{center}
{\Large MERIDIANKREIS.}

\rule{0.6\linewidth}{0.4pt}

Göttingische gelehrte Anzeigen. Nr. 91. S. 905..912. 5. Juni 1820.

\rule{0.6\linewidth}{0.4pt}
\end{center}

Der neue Meridiankreis von \textsc{Reichenbach}, den unsere Sternwarte im vergangenen Jahr erhalten hat, ist zwar schon seit geraumer Zeit aufgestellt und in Gebrauch. Ehe wir öffentlich darüber berichten, meinten wir jedoch, erst eine beträchtliche Reihe von Erfahrungen abwarten zu müssen, um danach umso sicherer beurteilen zu können, was dieses Instrument von einer bisher noch nicht verwendeten Bauart leistet.

Das Instrument dient zugleich als Mittagsfernrohr und als Höhenmesser. Für den ersten Zweck hat es mit den besten einfachen Mittagsfernrohren alle Einrichtungen gemeinsam. Sie bedürfen daher keiner ausführlichen Beschreibung. Das Fernrohr hat eine Brennweite von 5 Pariser Fuß und eine Öffnung von 4 Pariser Zoll. Die vier Okulare vergrößern 68-, 86-, 120- und 170-fach. Hofrat \textsc{Gauss} benutzt fast ausschließlich das stärkste. Das Fadennetz musste hier erneuert werden, weil einige Fäden schlaff geworden waren. Es besteht jetzt aus sieben senkrechten und zwei waagerechten Spinnenfäden. Sterne am Äquator durchlaufen die Zwischenräume zwischen den senkrechten Fäden in $\cdot 14$ Sekunden. Die waagerechten Fäden sind nur $7''6$ voneinander entfernt.

Die Achse ist 33 Pariser Zoll lang und trägt auf der einen Seite zwei konzentrische Kreise, deren vom Fernrohr abgewandte Flächen fast in derselben Ebene liegen. Der äußere Kreis sitzt \textit{fest} auf der Achse und dreht sich daher mit dem Fernrohr. Er trägt die Teilung, die unmittelbar in Schritten von 3 \SourceLeaf{435}
Minuten ausgeführt ist. Der innere Kreis (Alhidadenkreis) könnte sich ohne eine am Pfeiler angebrachte Hemmung um die Achse drehen. Die Hemmung erlaubt nur eine kleine, feine Bewegung zum Einstellen der am Kreis befestigten Wasserwaage. Auf diesem Alhidadenkreis liegen die vier Indizes jeweils $45^\circ$ von der Senkrechten entfernt. Ihre Nonien unterteilen jeden Hauptabschnitt in 90 Teile, also in Schritte von 2 Sekunden. Kleinere Teile lassen sich noch abschätzen. Der Durchmesser der Kreise an der Ablesestelle beträgt 35 Pariser Zoll. Die beiden Kreise berühren sich nicht und sind doch nur durch einen kaum merklichen Zwischenraum getrennt. Die Fläche des geteilten Kreises ragt äußerst wenig über die des Alhidadenkreises hinaus, weshalb die Ablesemikroskope absichtlich etwas schräg stehen. Diese Einrichtungen hat das Instrument mit anderen von \textsc{Reichenbach} gemeinsam. Drei Ausgleichsvorrichtungen für das ganze Instrument, den Alhidadenkreis und das Fernrohr sollen den Druck der Zapfen auf die Lagerschalen und des Alhidadenkreises auf die Achse sowie die Biegung des Rohrs aufheben. Ein Ausschlag von einem Pariser Zoll entspricht an der Hängelibelle, mit der die Achse waagerecht eingestellt wird, $22''$, an der Hauptlibelle $17''6$. Letztere dient dazu, den Alhidadenkreis stets in unveränderter Lage zu halten oder sehr kleine Verschiebungen zu messen und in Rechnung zu stellen. (Hofrat \textsc{Gauss} berührt die Korrekturschraube gewöhnlich nur, wenn die Verstellung mehr als $2''$ beträgt, was selten vorkommt.)

Das Instrument lässt sich mithilfe einer sehr zweckmäßig eingerichteten Maschine leicht und sicher umlegen. Bei der Höhenmessung dient das Umlegen zunächst dazu, den Kollimationsfehler zu ermitteln und die gemessenen Zenitdistanzen in absolute umzuwandeln. Es kann jedoch bedenklich erscheinen, den Kollimationsfehler über längere Zeit und bei beträchtlichen Temperaturänderungen als völlig unveränderlich anzusehen. Wenn man das Instrument nicht sehr häufig umlegt, ist es daher vorteilhafter, die Beobachtungen auf den Pol statt auf den Zenit zu beziehen. Der Ort des Pols auf dem Kreis lässt sich anhand von Zirkumpolarsternen umso leichter bestimmen, als die optische Leistung des Fernrohrs bewundernswert groß ist. Es steht darin dem Fernrohr am Passageinstrument von \textsc{Reichenbach} (d. Anz. 1819 [S. 422 d. B.]) kaum nach, obwohl dieses bei größerer Brennweite auch eine etwas größere Öffnung hat. Ein etwa noch vorhandener kleiner Unterschied wird dadurch ausgeglichen, dass man den Meridiankreis vorher aufs Genaueste einstellen kann und so jedes Mal seine ganze Aufmerksamkeit nur auf einen sehr kleinen Teil des Gesichtsfelds richten muss. Unter den von Hofrat \textsc{Gauss} am Tag beobachteten Sternen sind viele von der vierten Größe. 1 Draconis \textsc{Hevelii} und $\omega$ Cephei \textsc{Hevelii}, die ebenfalls am Tag und um die Mittagszeit beobachtet wurden, sind sogar nur von der fünften Größe. Dank der Einrichtung des Fadennetzes sind die Tagesbeobachtungen so feiner Pünktchen, sobald diese überhaupt sichtbar sind, ebenso zuverlässig wie die Beobachtungen der helleren Sterne.

\Needspace{40\baselineskip}
Als Beispiel für die ausgezeichnete Übereinstimmung, die sich mit diesem Instrument erreichen lässt, geben wir hier die Ergebnisse aller Bestimmungen des Orts des Pols auf dem Instrument von der Eröffnung des Beobachtungsbuchs bis zum ersten Umlegen. Sie stammen aus sämtlichen unmittelbar nacheinander beobachteten entgegengesetzten Kulminationen von Zirkumpolarsternen. Die Übereinstimmung der Ergebnisse für ein und denselben Stern beweist, wie genau sich die Beobachtungen anstellen lassen und dass der Kollimationsfehler in dieser Zeit unverändert blieb. Die ebenso große Übereinstimmung der Ergebnisse für verschiedene Sterne beweist dagegen die ausgezeichnete Qualität der Teilung von \textsc{Reichenbach}, die angesichts der mäßigen Größe des Kreises umso mehr Bewunderung verdient.

\Needspace{29\baselineskip}
\SourceLeaf{436}
\begin{center}
\begin{tabular}{l|l|c|c}
1820 & Name der Sterne & Größe & Ort des Pols \\
\hline
Febr. 26 & $\alpha$ Cephei & 3 & $321^\circ\,29'\,32''14$ \\
 & $\beta$ Cephei & 3 & $31.89$ \\
 & $\delta$ Draconis & 3 & $32.29$ \\
 & $\eta$ Cephei & 3.4 & $32.24$ \\
 & $\alpha$ Cephei & 3 & $31.16$ \\
 & $\beta$ Cephei & 3 & $31.65$ \\
27 & $\delta$ Draconis & 3 & $31.39$ \\
 & $\eta$ Cephei & 3.4 & $31.90$ \\
 & $\delta$ Draconis & 3 & $31.54$ \\
 & $\eta$ Cephei & 3.4 & $31.63$ \\
28 & $\gamma$ Cephei & 3 & $31.54$ \\
 & $\tau$ Cephei & 4 & $32.74$ \\
 & $\delta$ Cephei & 4 & $31.12$ \\
März 9 & $\delta$ Draconis & 3 & $31.18$ \\
 & $\beta$ Cephei & 3 & $32.18$ \\
 & $\tau$ Cephei & 4 & $31.42$ \\
 & $\iota$ Cephei & 4 & $31.52$ \\
10 & $\beta$ Cephei & 3 & $32.05$ \\
11 & 1 Dracon. Hev. & 5 & $31.97$ \\
18 & 1 Dracon. Hev. & 5 & $32.04$ \\
 & 3 Lacertae & 4 & $32.32$ \\
 & $\alpha$ Ursae maj. & 2 & $31.62$ \\
19 & $\alpha$ Ursae min. & 2 & $32.81$ \\
\end{tabular}
\end{center}

\par\indent Der Mittelwert beträgt $321^\circ\,29'\,31''84$. Keine einzige dieser 23 Bestimmungen weicht um eine volle Sekunde von ihm ab.

\par\indent Soweit man die Beobachtungen an diesem Instrument auf den Pol statt auf den Zenit bezieht, ist die genaueste Kenntnis der Polhöhe im Grunde von untergeordneter Bedeutung. Wo es auf eine einzelne Sekunde oder gar auf Teile von Sekunden ankommt, ist größte Vorsicht geboten. Da man auf Dauer keine \textit{so große} Konstanz des Kollimationsfehlers erwarten kann, wie sie die obige Übersicht zeigt, empfiehlt es sich für die genaueste Bestimmung der Polhöhe, nur Beobachtungen zu verbinden, die unmittelbar vor und unmittelbar nach dem Umlegen gemacht wurden. Hofrat \textsc{Gauss} wagt es daher nicht, seine bisherigen Beobachtungen für \textit{diese} Bestimmung zu verwenden. Nach dem ersten Umlegen herrschte ziemlich lange ungünstiges Wetter, und beim zweiten am 13. April ließ eine zufällige kleine Störung des Fadennetzes die \textit{vollständige} Gleichheit des Kollimationsfehlers etwas zweifelhaft erscheinen. Die Ergebnisse dieser Beobachtungen liegen im Übrigen sehr nahe beieinander und ebenso nahe an einem Ergebnis, das auf einem anderen Weg gewonnen wurde und das wir gleich anführen wollen. Verbindet man nämlich den Ort des Pols im ersten Zeitraum mit dem im zweiten ($38^\circ\,25'\,53''\!54$), so erhält man die Polhöhe $51^\circ\,31'\,49''\!15$. Verbindet man die 13 Zirkumpolarsternbeobachtungen unmittelbar vor dem zweiten Umlegen mit den unmittelbar danach beobachteten entgegengesetzten Kulminationen derselben Sterne, so erhält man $51^\circ\,31'\,47''\!86$.

Seit einigen Jahren hat bekanntlich ein Umstand die Aufmerksamkeit der Astronomen in hohem Maße auf sich gezogen. Obwohl es dabei nur um ein paar Sekunden geht, ist er sowohl für die Beobachtungskunst als auch für die vielfältigen astronomischen Ergebnisse von größter Bedeutung. Wir meinen die kleinen Unterschiede, die bei der Bestimmung von Sterndeklinationen, der Schiefe der Ekliptik und von Polhöhen mit verschiedenen, im Übrigen aber ganz ausgezeichneten Instrumenten aufgetreten sind. Ohne Zweifel gehen diese Unterschiede auf die \emph{Einwirkung der Schwerkraft} auf die verschiedenen Teile jedes Instruments zurück, auch wenn sich bis \SourceLeaf{437}
jetzt weder die Art der Einwirkung vollständig und mit Sicherheit nachweisen noch mit Bestimmtheit beurteilen lässt, welches Instrument das richtige und welches das unrichtige Ergebnis geliefert hat. Wie stark sich Metalle biegen, wissen wir im Grunde noch sehr wenig. Es scheint uns zu gewagt, bei irgendeinem Instrument, wie es auch gebaut sein mag, einen merklichen Einfluss der Schwerkraft auf seine Teile und damit auf die Beobachtungen für ausgeschlossen zu halten, ohne eine solche Behauptung durch hinreichende andere Belege zu begründen. Beim hiesigen Meridiankreis hat der große Instrumentenbauer zwar alles getan, um die Biegung des Fernrohrs durch eine ausgeklügelte Ausgleichsvorrichtung aufzuheben. Man könnte aber immer noch zweifeln, ob dadurch jede Flexion beseitigt oder in ihren Wirkungen völlig unmerklich geworden ist. Das einzige direkte Mittel zur Prüfung scheint darin zu bestehen, unmittelbare Beobachtungen eines Himmelskörpers mit Beobachtungen seines in einem künstlichen Horizont gesehenen Bildes zu verbinden. Solche Beobachtungen müssen selbstverständlich oft wiederholt werden, um diese heikle Frage zu klären. Hofrat \textsc{Gauss} hat mit solchen Beobachtungen bereits begonnen, indem er den Polarstern im Spiegel einer Wasseroberfläche beobachtete. Den glänzendsten Beweis für die erstaunliche optische Leistung des Fernrohrs liefert wohl der Umstand, dass sich auch die obere, am Tag eintretende Kulmination auf diese Weise recht gut beobachten ließ. Das Ergebnis der ersten \textit{vollständigen} Beobachtung dieser Art war folgendes:

\Needspace{22\baselineskip}
{\rightskip=0pt plus 3em \textit{Zenitdistanz des Polarsterns} (frei von Strahlenbrechung, aber einschließlich des Kulminationsfehlers)\par}

\begin{center}
1820 Mai 13.
\end{center}

\begin{center}
\begin{tabular}{l c l}
Untere Kulmination, direkt & . . . . . . . . . . & $319^\circ\,50'\,20''\!73$ \\
Untere Kulmination, reflektiert & . . . . . . . . . & $220\ \ 5\ \ 3.94$ \\
Obere Kulmination, direkt & . . . . . . . . . . & $323\ \ 8\ 41.51$ \\
Obere Kulmination, reflektiert & . . . . . . . . . & $216\ 46\ 44.31$ \\
\end{tabular}
\end{center}

\begin{center}
Hieraus folgt die \textit{wahre} Zenitdistanz
\end{center}

\begin{center}
\begin{tabular}{l c l}
in der unteren Kulmination & . . . . . . . . . . & $40\ \ 7\ 21.60$ \\
in der oberen Kulmination & . . . . . . . . . . & $36\ 49\ \ 1.40$ \\
\end{tabular}
\end{center}

\noindent und daraus ergibt sich bei einer Änderung der Deklination in 12 Stunden von $-0''10$ die Polhöhe $51^\circ\,31'\,48''\!45$ am Standort des Wassergefäßes oder $51^\circ\,31'\,48''\!40$ für den Mittelpunkt des Kreises.

\par\indent Da diese Bestimmung fast genau zwischen den beiden oben angeführten liegt, ist es bereits sehr wahrscheinlich, dass die Wirkung der Schwerkraft auf die Beobachtungen mit diesem Instrument entweder völlig unmerklich oder doch äußerst gering sein muss. Die Deklination des Polarsterns aus den obigen Beobachtungen ist um $0''42$ kleiner als nach \textsc{Bessels} Tafeln (sämtliche bisherigen Beobachtungen von Hofrat \textsc{Gauss} ergeben für diese Korrektur $=
-0''67$).

\Needspace{17\baselineskip}
Um auch ein Beispiel dafür zu geben, was das Instrument als Mittagsfernrohr leistet, führen wir hier einige Beobachtungen des Mars an, die Hofrat \textsc{Gauss} um die Zeit der Quadratur dieses Planeten gemacht hat.

\begin{center}
\begin{tabular}{l|c|c|c}
1820 & Mittlere Zeit & \shortstack{Rektasz. des\\1. Randes} & \shortstack{Dekl. des\\Mittelpunkts} \\
\hline
März 29 & $7^{\mathrm u}\,10^{\mathrm m}\,16^{\mathrm s}3$ & $114^\circ\,37'\,49''3$ & $24^\circ\,16'\,47''3$ N. \\
30 & $7\,7\,40.7$ & $114\,57\,57.7$ & $24\,12\,7.1$ \\
31 & $7\,5\,6.8$ & $115\,18\,30.1$ & $24\,7\,18.1$ \\
April 5 & $6\,52\,37.3$ & $117\,6\,19.9$ & $23\,42\,5.4$ \\
10 & $6\,40\,40.4$ & $119\,2\,17.5$ & $23\,14\,30.6$ \\
11 & $6\,38\,20\,3$ & $119\,26\,18.3$ & $23\,8\,41.0$ \\
12 & $6\,36\,1.6$ & $119\,50\,41.5$ & $23\,2\,48.0$ \\
13 & $6\,33\,43.7$ & $120\,15\,15.0$ & $22\,56\,44.8$ \\
26 & $6\,5\,16.8$ & $125\,56\,10.0$ & $21\,29\,12.8$ \\
\end{tabular}
\end{center}

\Needspace{18\baselineskip}
\SourceLeaf{438}
Herr von \textsc{Staudt} verglich diese Beobachtungen mit den Marstafeln von Herrn von \textsc{Lindenau}. Er widmet sich derzeit bei uns mit ausgezeichnetem Eifer und Erfolg dem Studium der mathematischen und astronomischen Wissenschaften. Es ergaben sich folgende

\begin{center}
Unterschiede

\begin{tabular}{ll|r|r}
 & & Rektasz. & Dekl. \\
\hline
März & 29 & $-0''7$ & $+5''1$ \\
 & 30 & $-0.5$ & $+3.9$ \\
 & 31 & $-2.1$ & $+6.1$ \\
April & 5 & $-1.3$ & $+3.9$ \\
 & 10 & $-4.9$ & $+2.9$ \\
 & 11 & $-1.9$ & $+3.6$ \\
 & 12 & $-4.8$ & $+1.9$ \\
 & 13 & $-1.8$ & $+4.2$ \\
 & 26 & $-1.5$ & $+1.5$ \\
\end{tabular}
\end{center}
